\emph{(i) The odd part.} The exchange $P_{12}$ fixes $u'$ and $e_{3}$, so $P_{12}\mathbb{1}P_{12} = \mathbb{1}$, $P_{12}\J P_{12} = \J$, $P_{12}EP_{12} = E$ (Lemma \claimref{neutrino.breaking-relation} for $\J$, $\Pi = \mathbb{1} - \J/3$ and $E$), and $u'\wedge s$ is $P_{12}$-invariant by Proposition \claimref{synthesis.residual-s2-fixes-120-family-form}(i). The shift $S$, $Se_{j} = e_{j+1}$ (indices mod $3$), has $S + S^{\mathsf T} = \J - \mathbb{1}$, even, and $S - S^{\mathsf T} = A$ with $Av = v\cdot(I_{e}u')$, $I_{e} = e_{1}e_{2}e_{3}$ the unit pseudoscalar of family space: the contraction of $v$ with the bivector dual to $u'$. A permutation matrix $P$ is orthogonal, so it extends to an automorphism of the algebra of family space sending $I_{e}$ to $\det P\,I_{e}$, and $P(v\cdot(I_{e}u')) = (Pv)\cdot(\det P\,I_{e}Pu') = \mathrm{sgn}(P)\,(Pv)\cdot(I_{e}u')$, so $PAP^{\mathsf T} = \mathrm{sgn}(P)A$: $A$ spans a copy of the sign representation, odd under every transposition and fixed by the $3$-cycles. The odd part $X \mapsto \tfrac12(X - P_{12}XP_{12})$ is linear; it kills the five even structures and sends $S$ and $S^{\mathsf T}$ to $\pm\tfrac12 A$, so the odd part of the complex span is $\C A$, and without $S$ and $S^{\mathsf T}$ (which enter only through the circulant reading of (ii)) it is zero. The sign representation occurs once in $M_{3}$: by characters, $M_{3} = V\otimes V$ with $V = \mathbf{1}\oplus\mathbf{2}$ the permutation representation, and $V\otimes V = 2\cdot\mathbf{1}\oplus\mathbf{sgn}\oplus 3\cdot\mathbf{2}$, since $\mathbf{2}\otimes\mathbf{2} = \mathbf{1}\oplus\mathbf{sgn}\oplus\mathbf{2}$; the check computes the space $\{X : PXP^{\mathsf T} = \mathrm{sgn}(P)X\ \forall P\}$ as an exact kernel and finds dimension one. It is antisymmetric (the symmetric matrices, $\mathrm{Sym}^{2}V = 2\cdot\mathbf{1}\oplus 2\cdot\mathbf{2}$, contain no sign representation).

\emph{(ii) The circulant reading.} From $e^{I_{4}\delta}S + e^{-I_{4}\delta}S^{\mathsf T} = \cos\delta\,(S + S^{\mathsf T}) + I_{4}\sin\delta\,(S - S^{\mathsf T})$ and (i), $C = \mathbb{1} + \tfrac{\lambda}{2}(e^{I_{4}\delta}S + e^{-I_{4}\delta}S^{\mathsf T}) = (1 - \tfrac{\lambda}{2}\cos\delta)\mathbb{1} + \tfrac{\lambda}{2}\cos\delta\,\J + I_{4}\tfrac{\lambda}{2}\sin\delta\,A$. On $f_{k} = (1, \zeta^{k}, \zeta^{2k})$, $(Sf_{k})_{i} = (f_{k})_{i - 1} = \zeta^{-k}(f_{k})_{i}$, so $Cf_{k} = (1 + \lambda\cos(\delta - 2\pi k/3))f_{k}$: the components of the projected vector $p$ of Proposition \claimref{synthesis.coherent-face-vector-real-slices}(iii), up to the overall $|\kappa|$ and the relabelling $k \mapsto -k$. The check verifies the matrix identity and the three eigenvector equations exactly in $\Q(\sqrt3)(i)$ at $\cos\delta = \tfrac35$, $\sin\delta = \tfrac45$, for $\lambda = 1$ and $\tfrac32$. So, in the reading in which the face family lines are the Fourier lines of a flavour basis, the $S_{3}$-invariant part of the square-root mass matrix is $\mathrm{span}\{\mathbb{1}, \J\}$ and the channel phase is the coefficient, $I_{4}\tfrac{\lambda}{2}\sin\delta$, of the sign line $A$; at $\delta = 0$ the light pair is degenerate (the skeleton), and $\delta$ breaks $S_{3}$ to $\mathbb{Z}_{3}$. That reading is a choice of basis, not a premise of the store.

\emph{(iii) Moduli.} Write $x = \hat{b}_{e}^{\mathsf T}X\hat{b}_{o}$ and $y = \hat{b}_{o}^{\mathsf T}X\hat{b}_{e}$, the entries of the compression of Proposition \claimref{synthesis.solar-angle-light-plane-sum-rule}(i) (the diagonal entries of a $P_{12}$-odd $X$ vanish, as there, the derivation being linear and valid over $\C$). Since $\hat{b}_{e}, \hat{b}_{o}$ are real, $X^{\mathsf T} = X$ gives $y = x$, $X^{\mathsf T} = -X$ gives $y = -x$, and $X^{\dagger} = X$ gives $y = \overline{\hat{b}_{e}^{\mathsf T}X\hat{b}_{o}} = \bar{x}$; in all three $|x| = |y|$. For $X = \gamma A$: $A\hat{b}_{o} = u'\times(1,-1,0)/\sqrt2 = (1,1,-2)/\sqrt2 = \sqrt3\,\hat{b}_{e}$ and $A\hat{b}_{e} = -\sqrt3\,\hat{b}_{o}$, so $x = \sqrt3\gamma = -y$ (the check works with the unnormalised $b_{e} = (1,1,-2)$, $b_{o} = (1,-1,0)$, where the entries are $\pm 6\gamma$). By contrast $AE = (Ae_{3})e_{3}^{\mathsf T} = (1,-1,0)e_{3}^{\mathsf T}$ is $P_{12}$-odd with $x = 0 \ne y$: maximally lopsided.

\emph{(iv) The light angle.} Let $B = \bigl[\begin{smallmatrix} k & x\\ y & 0\end{smallmatrix}\bigr]$, $k > 0$, $x, y \in \C$. Then $BB^{\dagger} = \bigl[\begin{smallmatrix} k^{2} + |x|^{2} & k\bar{y}\\ ky & |y|^{2}\end{smallmatrix}\bigr]$, with characteristic polynomial $q(s) = s^{2} - (k^{2} + |x|^{2} + |y|^{2})s + |x|^{2}|y|^{2}$, and $u = (k\bar{y},\ s - k^{2} - |x|^{2})$ satisfies $(BB^{\dagger} - s)u = (0, -q(s))$, so for the light root $s = \sigma_{2}^{2}$ it is the light left singular vector, with $|u_{1}/u_{2}| = k|y|/(\sigma_{1}^{2} - |y|^{2})$, using $\sigma_{1}^{2} + \sigma_{2}^{2} = k^{2} + |x|^{2} + |y|^{2}$. Only the moduli enter. With $|x| = |y| = r$: $\sigma_{1}\sigma_{2} = r^{2}$ and $k^{2} = (\sigma_{1} - \sigma_{2})^{2}$, so $\tan^{2}\theta = k^{2}r^{2}/(\sigma_{1}^{2} - r^{2})^{2} = (\sigma_{1} - \sigma_{2})^{2}\sigma_{1}\sigma_{2}/(\sigma_{1}^{2}(\sigma_{1} - \sigma_{2})^{2}) = \sigma_{2}/\sigma_{1}$; equivalently $k^{2}r^{4} = s(k^{2} + r^{2} - s)^{2}$ modulo $q$, which the check verifies by exact polynomial division at four values of $(k, r^{2})$, together with the eigenvector relation at two complex pairs of equal modulus. So, as in the real case, the light charged state makes the angle $\theta_{e} = \arctan\sqrt{m_{e}/m_{\mu}}$ with $\hat{b}_{o}$ in modulus, now with complex components. For a unit vector $e$ with $|e\cdot\hat{b}_{e}| = \sin\theta_{e}$, $|e\cdot\hat{b}_{o}| = \cos\theta_{e}$ and a real unit vector $\nu = \sin\theta_{\nu}\hat{b}_{e} + \cos\theta_{\nu}\hat{b}_{o}$ (the real neutrino stage of the sum rule, $\tan\theta_{\nu} = (m_{1}/m_{2})^{1/4}$), the triangle inequality gives $|\langle e, \nu\rangle| \ge \cos\theta_{e}\cos\theta_{\nu} - \sin\theta_{e}\sin\theta_{\nu} = \cos(\theta_{e} + \theta_{\nu})$, which the check confirms in the form $|1 + t_{e}t_{\nu}w|^{2} \ge (1 - t_{e}t_{\nu})^{2}$, $|w| = 1$, at exact samples. Hence $|U_{e1}| \ge \cos(\theta_{e} + \theta_{\nu})$ and the sum rule $\Delta \le \arctan\sqrt{m_{e}/m_{\mu}} + \arctan(m_{1}/m_{2})^{1/4}$ of Proposition \claimref{synthesis.solar-angle-light-plane-sum-rule}(iii) holds with the charged stage complex, $|x_{e}| = |y_{e}|$; under the inequality of Premise (H) (\claimref{neutrino.h-second-stage-hierarchical}), taken as a hypothesis of this clause as in that proposition, the stored numbers apply: $\Delta < 21.6^{\circ}$ and $|U_{e1}|^{2} > 0.865$, against $|U_{e1}|^{2} \le 0.711$.

\emph{(v) Consequence.} By (i) a second stage built from the series' structures has $X_{e} = \gamma_{e}A$ and $X_{\nu} = \gamma_{\nu}A$; by (iii) both are one type in modulus; by (iv) with $\gamma_{\nu}$ real (the stored setting) the measured solar angle is excluded under (H). A complex neutrino coefficient is outside this proof: a complex Dirac stage enters $m_{\nu} = \mD M_{R}^{-1}\mD^{\mathsf T}$ through the transpose, and the light neutrino angle is no longer a function of the moduli. In exploration (floating point, not part of the check) the largest light neutrino angle over complex $\gamma_{\nu}$ with $m_{1}/m_{2} \le 0.01$ was about $19.3^{\circ}$, against $17.5^{\circ}$ for real $\gamma_{\nu}$, so with the charged $4.0^{\circ}$ the total stays near $23^{\circ}$, far below the measured $\Delta \ge 32.5^{\circ}$; this is recorded as a numerical observation. What the solar angle needs, $|y| \ne |x|$ in some sector, is an $S_{2}$-odd matrix with a symmetric part, that is a component of a standard representation, together with an antisymmetric part at a relative phase that is not Hermitian; no linear combination of the series' structures supplies it, and products of them can (for instance $AE$, which is a product and not a coupling of the store's textures).
